\chapter{מהוכחת היתכנות לסביבת Production}
\section{הפער השקט בין Demo ל-Production}
ההבדל בין Demo מרשים למערכת Production אינו נמדד באיכות פלט בודד, אלא בהתנהגות המערכת תחת שינוי, כשל ועומס. ב-Demo כל Agent מקבל קלט נקי, ה-Crew מריץ Process יחיד בנתיב מאושר מראש, וה-Context עובר בין השלבים בדיוק כמצופה. ברגע שמערכת התזמור פוגשת קלטים אמיתיים נחשף הפער השקט שאיש אינו רואה במצגת: מה קורה כש-Task נכשל באמצע השרשרת, כשמודל מחזיר LaTeX שבור או JSON שאינו עובר validation, או כששירות חיצוני קורס. מערכת Production מחייבת recovery מובנה — ניסיון חוזר עם backoff, נתיב fallback, ומנגנון שמבודד כשל של Agent בודד כדי שלא יפיל את ה-Crew כולו. לצד ההתאוששות נדרשת monitoring רציפה: לא לוג יחיד בסוף הריצה, אלא observability בכל שלב, עם מדדי זמן, שיעורי כשל, ועקיבות מלאה אחר ה-Context העובר בין ה-Tasks. שכבת ה-Harness הדטרמיניסטית קריטית כאן, משום שהיא מאפשרת לשחזר ריצה, להשוות פלטים, ולזהות רגרסיה כשמודל או prompt משתנים. בלי control מפורש על השינוי — גרסאות נעולות של מודלים, בדיקות שמריצות תרחישים מוכרים לפני פריסה, ושערי validation החוסמים פלט לא תקין — המערכת תתפקד היטב היום ותיכשל באופן בלתי צפוי מחר. המעבר ל-Production הוא במידה רבה מעבר מאופטימיזציה למקרה הטוב לעיצוב למקרה הגרוע, שבו כל רכיב מתוכנן מראש להיכשל בחן ולהתאושש באופן נשלט.

\section{אבטחת Agents ומודעות לעלות}
כשמערכת תזמור עוברת ל-Production, כל Agent הופך הן למשטח תקיפה פוטנציאלי והן למקור הוצאה כספית, ושני הממדים הללו דורשים משמעת הנדסית. עקרון ה-least privilege הוא הבסיס: כל Agent מקבל אך ורק את הכלים, ההרשאות והגישה לנתונים הדרושים למשימתו, ולא יותר מכך. Agent שתפקידו לעצב LaTeX אינו זקוק לגישה לרשת או למפתחות API חיצוניים, ו-Agent שמסכם נתונים אינו צריך הרשאת כתיבה למאגר. הפרדה זו מצמצמת את הנזק כאשר prompt injection או קלט עוין מצליחים לחדור. מעליה נדרשים guardrails — שכבות validation הבודקות הן את הקלט הנכנס והן את הפלט היוצא, מסננות הוראות זדוניות המוטמעות ב-Context, ומוודאות שהפעולות שה-Agent מבקש לבצע נמצאות בתחום המותר. במקביל, מודעות לעלות אינה מותרות אלא חלק מהארכיטקטורה: כל קריאה למודל צורכת tokens שעולים כסף, ושרשרת Tasks ארוכה עם Context נפוח עלולה לייקר ריצה בודדת פי כמה. משמעת של token-cost מחייבת מדידה לכל Task, תקצוב מרבי לריצה, בחירה מושכלת בין מודל יקר וחזק למודל זול וקל למשימות פשוטות, וצמצום ה-Context המועבר לרלוונטי בלבד. ניטור עלות לצד observability מאפשר לזהות Agent שנכנס ללולאה יקרה או prompt שמנפח צריכה. אבטחה ועלות הם שני צדדים של אותה משמעת: שליטה הדוקה במה שכל Agent רשאי לעשות ובמה שכל Agent מבזבז.

\section{מחסנית מודולרית והכיוון הבא}
מערכת Production בוגרת נשענת על מחסנית מודולרית ומנוטרת, שבה כל רכיב — הגדרת Agent, Task, בחירת Process, שכבת ה-Harness ושערי ה-validation — ניתן להחלפה, לבדיקה ולניטור בנפרד. מודולריות זו אינה נוחות הנדסית בלבד אלא תנאי להתפתחות: כשהגבולות בין הרכיבים ברורים ו-observability מכסה כל שלב, ניתן לשדרג רכיב אחד מבלי לסכן את השלם, ולהריץ ניסויים מבוקרים. מנקודה זו הארכיטקטורה ממשיכה בכמה כיוונים טבעיים. הראשון הוא hierarchical crews — במקום Crew שטוח שבו כל Agent שווה, נוצר Agent מנהל המתזמר Agents כפופים, מפרק בעיה מורכבת לתת-משימות, ומקצה אותן באופן דינמי. מבנה היררכי כזה מאפשר להתמודד עם משימות פתוחות שאינן ניתנות לתכנון מראש ב-Process קווי. הכיוון השני הוא שילוב RAG, שבו ה-Agents שואבים ידע עדכני ומבוסס ממאגר חיצוני במקום להישען על משקלי המודל בלבד, וכך משפרים דיוק ומצמצמים hallucination בתחומי ידע ספציפיים. הכיוון השלישי, ואולי החשוב מכולם לאורך זמן, הוא evaluation שיטתי: מסגרת מדידה שמריצה את ה-Crew על אוסף תרחישים מייצגים, מציינת את איכות הפלטים מול קריטריונים מוגדרים, ומזהה רגרסיה לפני פריסה. בלי evaluation אובייקטיבי כל שיפור נותר ניחוש. שילוב של מחסנית מודולרית, hierarchical crews, RAG ו-evaluation הוא הנתיב שבו מערכת תזמור הופכת מכלי חד-פעמי לפלטפורמה המתפתחת, נמדדת ומשתפרת באופן שיטתי.

לקריאה נוספת בנושאי פרק זה, ראו \cite{latex_project}.

